% !TEX root = userguide.tex

\def\headdate{4th June 2020}
\label{sec:small_lambda}

\section{Handling small relaxation times in multi-mode viscoelastic fluid flow}

(This section has been contributed by Nick Jaensson of the TU/e).\medskip

Most time integration schemes in Secs.~\ref{sec:explicitstress} and
\ref{sec:implicitstress} have some form of implicit treatment, however the
relaxation terms of the constitutive model are always explicit. Multi-mode
models of real polymers usually have a broad spectrum with also very small
relaxation times. This would limit the time step enormously, since it must be
smaller than the smallest relaxation time. A possible solution is to use fully
implicit schemes, but this is expensive since it requires iteration and the
system matrix becomes larger. An ad hoc approach is to 
replace the small relaxation time modes by a viscous component, but this is not
very convenient and also fully ignores the non-linear behavior of a mode.

In this section we propose an approach which gradually reduces the smallest
relaxation time modes to Newtonian behavior. The cut-off depends on the time
step taken.

\subsection{Increasing the time step limit for small relaxation times}

The approach taken here is to add
the term $\ten c/\lambda$ to both sides of the constitutive 
equation, and to use an implicit time discretization for this term on the LHS, 
while using an explicit discretization of this term on the RHS. Note, that
$\lambda$ is the viscoelastic relaxation time constant for the mode at hand.
This approach assures that for sufficiently small time step the mode is fully
resolved. For large time steps and small ``mode Deborah number'', Newtonian
behavior is approached. For larger ``mode Deborah number'' it depends on the
behavior of the non-linear terms whether stable steady state
results are obtained. In any case,
transient behavior is incorrect for time steps large compared to the
relaxation time.

For ease of reading, we begin by repeating the viscoelastic constitutive 
equation:
\begin{equation}
  \pderiv{\ten c}{t}
  +\vek u\cdot\nabla\ten c-\ten L\cdot\ten c-\ten c\cdot\ten L^T
   +\ten f(\ten c)=\ten 0 \label{eq:const_ve}
\end{equation}
Using the definition $\ten g(\ten L,\ten c)=\ten L\cdot\ten c
+\ten c\cdot\ten L^T-\ten f(\ten c)$, Eq.~\eqref{eq:const_ve} is written as
\begin{equation}
  \pderiv{\ten c}{t}
  +\vek u\cdot\nabla\ten c=\ten g(\ten L,\ten c) \label{eq:const_ve2}
\end{equation}
It is known that the Crank-Nicolson scheme can yield solutions that
oscillate around the real solution when the time step is chosen large.
Therefore, this approach is only implemented for 
the first-order semi-implicit Euler scheme and the second-order
semi-implicit Gear scheme, which both do not suffer from these
oscillations.

\subsubsection{Semi-implicit Euler}
Adding the $\ten c/\lambda$ term to both sides of Eq.~\eqref{eq:const_ve2} 
and using a semi-implicit Euler scheme for the discretization yields:
\begin{equation}
  \frac{\ten c^{n+1}-\ten c^n}{\Delta t}+\vek u^{n+1}\cdot\nabla\ten c^{n+1}+
  \frac{\ten c^{n+1}}{\lambda}=
\ten g(\ten L^{n+1}, \ten c^n)+
  \frac{\ten c^{n}}{\lambda}
  \label{eq:euler_c_over_lambda}
\end{equation}
The effectiveness of the additional terms can be seen by substituting the 
UCM/Oldroyd-B model $\ten f(\ten c)=(\ten c - \ten I)/\lambda$ into
Eq.~\eqref{eq:euler_c_over_lambda}:
\begin{equation}
  \frac{\ten c^{n+1}-\ten c^n}{\Delta t}+\vek u^{n+1}\cdot\nabla\ten c^{n+1}+
  \frac{\ten c^{n+1}}{\lambda}=
\ten L^{n+1}\cdot\ten c^n + \ten c^n \cdot 
  {\ten{L}^{n+1}}^T +
\frac{\ten c^{n}}{\lambda}-\frac{1}{\lambda}(\ten c^n-\ten I)
\end{equation}
which can be rewritten to
\begin{equation}
  \frac{\ten c^{n+1}-\ten c^n}{\Delta t}+\vek u^{n+1}\cdot\nabla\ten c^{n+1}+
  \frac{1}{\lambda}(\ten c^{n+1}-\ten I)=
  \ten L^{n+1}\cdot\ten c^n + \ten c^n \cdot {\ten{L}^{n+1}}^T
\end{equation}
Note that the third term on the LHS equals $\ten f(\ten c^{n+1})$. So, for 
the UCM/Oldroyd-B model, 
adding the $\ten c/\lambda$ implicitly/explicitly has the net effect 
of treating the $\ten f(\ten c)$ term implicitly. 
For the implementation in 
\tfem, Eq.~\eqref{eq:euler_c_over_lambda} is rewritten to
\begin{equation}
  \left(\frac{1}{\Delta t}+\frac{1}{\lambda}\right)\ten c^{n+1}+
  \vek u^{n+1}\cdot\nabla\ten c^{n+1}=
     \frac{\ten c^n}{\Delta t} +
     \ten g(\ten L^{n+1},\ten c^n) + \frac{\ten c^n}{\lambda}
  \label{eq:Eulerimp_add}
\end{equation}
This form is used for the implementation, because the $\ten c^n/\lambda$ term 
can be immediately combined with the $\ten g(\ten L^{n+1},\ten c^n)$ term.

\subsubsection{Semi-implicit Gear with prediction}
Adding the $\ten c/\lambda$ term to both sides of Eq.~\eqref{eq:const_ve2} 
and using a second order semi-implicit Gear scheme for the discretization
yields
(note that the $\ten c/\lambda$ term on the RHS is based on the 
prediction $\hat{\ten c}^{n+1}$):
\begin{equation}
  \frac{\frac32 \ten c^{n+1}-2\ten c^n+\frac12 \ten c^{n-1}}{\Delta t}+
  \vek u^{n+1}\cdot\nabla\ten c^{n+1}+\frac{\ten c^{n+1}}{\lambda}=
\ten g(\ten L^{n+1},\hat{\ten c}^{n+1})+
  \frac{\hat{\ten c}^{n+1}}{\lambda}
  \label{eq:gear_c_over_lambda}
\end{equation}
where $\hat{\ten c}^{n+1}=2\ten c^n-\ten c^{n-1}$ is a prediction for the 
conformation tensor. In a similar way as was shown for the semi-implicit 
Euler scheme, the additional terms have the net effect of treating the
$\ten f(\ten c)=(\ten c - \ten I)/\lambda$ term implicitly for the
UCM/Oldroyd-B model.
For the implementation in 
\tfem, Eq.~\eqref{eq:gear_c_over_lambda} is rewritten to
\begin{equation}
 \left(\frac{3}{2\Delta t}+\frac{1}{\lambda}\right)\ten c^{n+1} +
   {\vek u}^{n+1}\cdot\nabla\ten c^{n+1}=
  \frac{2\ten c^{n}-\frac12\ten c^{n-1}}{\Delta t} +
    \ten g(\ten L^{n+1},\hat{\ten c}^{n+1}) +
       \frac{\hat{\ten c}^{n+1}}{\lambda}  \label{eq:gear_implementation}
\end{equation}
This form is used for the implementation, where the $\hat{\ten 
c}^{n+1}/\lambda$ term can be immediately combined with the $\ten 
g(\ten L^{n+1},\hat{\ten c}^{n+1})$ term.

\subsubsection{Semi-implicit Gear (with explicit Karniadakis)}
Adding the $\ten c/\lambda$ term to both sides of Eq.~\eqref{eq:const_ve2} 
and using a second order semi-implicit Gear scheme (with explicit Karniadakis)
for the discretization yields
\begin{multline}
  \frac{\frac32 \ten c^{n+1}-2\ten c^n+\frac12 \ten c^{n-1}}{\Delta t}+
  \vek u^{n+1}\cdot\nabla\ten c^{n+1}+\frac{\ten c^{n+1}}{\lambda}=\\
2\big[\ten g(\ten L^n,{\ten c}^{n})+\frac{{\ten c}^{n}}{\lambda}\big]
-\big[\ten g(\ten L^{n-1},{\ten c}^{n-1})+
  \frac{\ten c^{n-1}}{\lambda}\big]
  \label{eq:geark_c_over_lambda}
\end{multline}
In a similar way as was shown for the semi-implicit 
Euler scheme, the additional terms have the net effect of treating the
$\ten f(\ten c)=(\ten c - \ten I)/\lambda$ term implicitly for the
UCM/Oldroyd-B model.
For the implementation in 
\tfem, Eq.~\eqref{eq:gear_c_over_lambda} is rewritten to
\begin{multline}
 \left(\frac{3}{2\Delta t}+\frac{1}{\lambda}\right)\ten c^{n+1} +
   {\vek u}^{n+1}\cdot\nabla\ten c^{n+1}=
  \frac{2\ten c^{n}-\frac12\ten c^{n-1}}{\Delta t} \\+
2\big[\ten g(\ten L^n,{\ten c}^{n})+\frac{{\ten c}^{n}}{\lambda}\big]
-\big[\ten g(\ten L^{n-1},{\ten c}^{n-1})+
  \frac{\ten c^{n-1}}{\lambda}\big]
  \label{eq:geark_implementation}
\end{multline}
This form is used for the implementation, where $\ten c/\lambda$
can be immediately combined with $\ten 
g(\ten L,\ten c)$ for the terms at $t^n$ and $t^{n-1}$.

Some remarks on using these adapted schemes:
\begin{itemize}
   \item For all schemes the factor in front of $\ten c^{n+1}$ 
         now depends on the relaxation time $\lambda$, and therefore the 
         system matrix needs to be rebuild for every mode if multiple 
         modes are used.
   \item For $\Delta t$ several times smaller than $\lambda$ the mode will be
         time-resolved. 
   \item For $\Delta t\gg\lambda$ the time-dependent term will 
         become very small compared to the relaxation terms. Therefore, the
         mode goes into a `steady state'', which is solved iteratively
         in the time stepping process. It now depends on the local Deborah
         number $\text{De}=\lambda |\ten D| \ll 1$ what will happen:
         \begin{itemize}
            
   \item If the $\text{De}\ll 1$, Newtonian    
         behavior will be recovered for $\Delta t/\lambda \gg 1$.
         This can be seen by expanding $\ten c$ in 
         a power series in $\lambda \ten D$, which yields
         \begin{equation}
           \ten c = \ten I + 2 \lambda \ten D +  
          \ten O\left( |\lambda \ten D|^2 \right) \label{eq:small_de}
         \end{equation}
         With the stress expression $\tau=G(\ten c-\ten I)$, Newtonian
         behavior with viscosity $\eta=G\lambda$ is obtained.
   \item For larger local Deborah numbers non-linear terms become important and
         whether the approach obtains stable results must be found by numerical
         experiments. At the time of writing not much experience is available.
         \end{itemize}

   \item For small Deborah numbers it follows from Eq.~\eqref{eq:small_de} that
         $|\ten c - \ten I| \approx |2\lambda \ten D| \ll 1$. Therefore, 
         $\ten f(\ten c)\approx(\ten c-\ten I)/\lambda$ and the Oldroyd/Maxwell
         behavior is obtained 
         and the approach is also valid if $\ten f(\ten c)$ is non-linear, but
         local Deborah is small.

\end{itemize}

For the stress-implicit method of Sec.~\ref{sec:implicitstress},
terms are also added to the prediction of
the conformation tensor that is used in the momentum equation in the stress 
implicit method. The stress tensor in the momentum balance is predicted using 
the semi-implicit Euler scheme as given in~\eqref{eq:Eulerimp_add}. 
In a similar way as in \ref{sec:implicitstress}, we can write the stress 
tensor as
\begin{equation}
  \begin{split}
   \ten \tau(\ten c^*{}^{n+1})&=
 \ten \tau\left(\frac{
          \Big(\Delta t(-\vek u^{n+1}\cdot\nabla\ten c^n+
          \ten L^{n+1}\cdot\ten c^n+\ten c^n\cdot\ten L^{n+1}{}^T)
          +\Big(1+\frac{\Delta t}{\lambda}\Big)\ten c^{n}-
           \Delta t\ten f(\ten c^n)\Big)}{\Big(1+\frac{\Delta t}{\lambda}\Big)}
      \right)\\
   &= \ten \tau\left(
           \frac{\Delta t}{1+\frac{\Delta t}{\lambda}}(-\vek u^{n+1}
           \cdot\nabla\ten c^n+
           \ten L^{n+1}\cdot\ten c^n+\ten c^n\cdot\ten L^{n+1}{}^T)
           +\ten h^*(\ten c^n, \Delta t, \lambda )
      \right)\\ 
  \end{split}
\end{equation}
where 
\begin{equation}
 \ten h^*(\ten c, \Delta t, \lambda)= \ten c-\Big(\frac{\Delta t}
  {1+\frac{\Delta t}{\lambda}}\Big)\ten f(\ten c)
\end{equation}
In case of a linear stress expression:
\begin{equation}
     \ten\tau=G(\ten c - \ten 1)
\end{equation}
we find
\begin{equation}
   \ten \tau(\ten c^*{}^{n+1})=
           \frac{G \Delta t}{1+\frac{\Delta t}{\lambda}}(-\vek u^{n+1}
           \cdot\nabla\ten c^n+
           \ten L^{n+1}\cdot\ten c^n+\ten c^n\cdot\ten L^{n+1}{}^T)
           +G\ten h^*(\ten c, \Delta t, \lambda )
            -G\ten 1
\end{equation}
Substituting this stress expression into the momentum equation leads to 
(for simplicity we set $\eta_s=0$):
\begin{align}
      \nabla p^{n+1} 
  -\alpha\nabla\cdot\big(\nabla\vek u^{n+1} -\ten G^T{}^{n+1} \big)
  -\nabla\cdot\big(\frac{G\Delta t}{1+\frac{\Delta t}{\lambda}}
   (&-\vek u^{n+1}\cdot\nabla\ten c^n+
        \ten L^{n+1}\cdot\ten c^n+
      \ten c^n\cdot\ten L^{n+1}{}^T)\big)\notag\\
&=\nabla\cdot(G \ten h^*(\ten c^n,\Delta t, \lambda) -G\ten 1)
        + \rho\vek b\\
     \nabla\cdot\vek u^{n+1}&=0\\
   -\nabla\vek u^{n+1} + \ten G^T{}^{n+1} &= \ten 0
\end{align}
\newpage
\noindent
Remark:
\begin{itemize}
  \item For local $De\ll1$ and $\Delta t/\lambda\gg1$ the additional implicit
terms tend to $\nabla\cdot(2G\lambda\ten D)$, i.e.\ the
additional viscosity is $\eta=G\lambda$.
\end{itemize}
If the log-conformation representation is used ($\ten s = \log \ten c$, see
Eq.~\eqref{eq:LCR})), the 
first-order semi-implicit Euler scheme given by
Eq.~\eqref{eq:Eulerimp_add} is written as
\begin{equation}
  \left(\frac{1}{\Delta t}+\frac{1}{\lambda}\right)\ten s^{n+1}+
  \vek u^{n+1}\cdot\nabla\ten s^{n+1}=
     \frac{\ten s^n}{\Delta t}+
      \ten h(\ten L^{n+1},\ten s^n) +
     \frac{\ten s^n}{\lambda}
\end{equation}
The second-order semi-implicit Gear scheme given by 
Eq.\eqref{eq:gear_implementation} changes in a similar way
\begin{equation}
 \left(\frac{3}{2\Delta t}+\frac{1}{\lambda}\right)\ten s^{n+1} +
   {\vek u}^{n+1}\cdot\nabla\ten s^{n+1}=
  \frac{2\ten s^{n}-\frac12\ten s^{n-1}}{\Delta t} +
    \ten h(\ten L^{n+1},\hat{\ten s}^{n+1}) +
       \frac{\hat{\ten s}^{n+1}}{\lambda}
\end{equation}
where $\hat{\ten s}^{n+1}=2\ten s^n-\ten s^{n-1}$ is a prediction for the 
conformation tensor in the logarithmic representation.

\subsection{Axi-symmetric problem: flow around a sphere in
 a cylindrical tube filled with a multi-mode model: handling 
 small relaxation times}
\label{sec:sphere_small_lambda}

The first example file is \texttt{sphere23.f90} which simulates a sphere in a
viscoelastic fluid with UCM two modes. Modes can have very short
relaxation times.

The example can be run by
\begin{quote}
   \texttt{sphere23 < input\_sphere23.txt}
\end{quote}
which might take a while. Change the parameters in the file
\texttt{input\_sphere23.txt}, if necessary.

Some notes on the source:
\begin{itemize}
   \item 
     \texttt{coefficients\%i(57)=1} is used to indicate that the
     $\ten c/\lambda$ term is added to both sides of the
     constitutive equation.
   \item 
     If these adapted time schemes are used, the system matrix to solve the
     constitutive equation is different for each mode. Therefore, 
     \texttt{build\_system} has to be included in the loop over the modes. 
     Furthermore, the current mode \texttt{m} has to be passed on to the 
     element routine.
     In order to do this, \texttt{coefficients\%i(26)=m} and 
     \texttt{coefficients\%i(27)=m} are used, and 
     \texttt{msysvector=rhsc(:,m)} is used in \texttt{build\_system} 
     (instead of \texttt{m2sysvector=rhsc} in most of the other 
     viscoelastic examples). Note, that any \texttt{build\_system\_constraint}
     should use \texttt{msysvector=rhsc(:,m)} as well.
   \item
     In the momentum balance, all modes are needed, so both 
     \texttt{coefficients\%i(26)} and \texttt{coefficients\%i(27)} 
     are set to zero after solving the constitutive equation for all modes.
   \item
     To show the convergence to Newtonian behavior for all modes having very
     small relaxation times, the output is the 
     difference in velocity of the current time step and the previous 
     time step. This is off by default and can be set by
     \texttt{printstepvel=.true.} in the source.
   \item The postprocessing \texttt{sphere\_post} can be used, but the mode
     which is output must be set with \texttt{outpmode}.
\end{itemize}
The second example file is \texttt{sphere24.f90} and is similar to
\texttt{sphere23.f90}, but uses a four mode Giesekus with data from literature.
